\documentclass[11pt]{article}
\usepackage[margin=1in]{geometry}

% Core packages
\usepackage{amsmath,amssymb}
\usepackage{array}
\usepackage{booktabs}
\usepackage{enumitem}
\usepackage{hyperref}
\usepackage{xcolor}

% Paragraphs
\setlength{\parindent}{0pt}
\setlength{\parskip}{0.8\baselineskip}
\emergencystretch=2em

% Lists
\setlist[itemize]{topsep=0.25\baselineskip,itemsep=0.25\baselineskip}
\setlist[enumerate]{topsep=0.25\baselineskip,itemsep=0.25\baselineskip}

% Table helpers
\newcolumntype{L}[1]{>{\raggedright\arraybackslash}p{#1}}

\title{Architecture Decision Record: ADR-003\\
\large Implementing AI-to-AI Communication over Zenoh}
\author{Abel-Raul Mazilu}
\date{13-05-2026}

\begin{document}
\maketitle

\begin{center}
\begin{tabular}{rl}
\textbf{Status:} & Accepted \\
\textbf{Project:} & DVerse Communication Platform 2.0 \\
\end{tabular}
\end{center}

\section{Context}

The DVerse platform is moving toward a multi-agent architecture in which AI agents act as autonomous participants. These agents need a decoupled communication mechanism that does not depend on direct service-to-service calls or the chat application transport layer.

Zenoh has already been selected as the messaging layer for the platform. Its topic-based publish-subscribe model is suitable for asynchronous agent communication and future secured transport requirements.

\section{Decision}

Sprint 4 will implement AI-agent-to-AI-agent communication over Zenoh. Agents will publish and subscribe to structured Zenoh topics as first-class participants in the system.

This implementation will include the following capabilities:

\begin{itemize}
    \item Topic conventions for agent discovery, commands, responses, status updates, and errors.
    \item Structured message schemas for payloads exchanged between agents.
    \item Publish and subscribe flows that allow agents to communicate without direct coupling.
    \item Message correlation identifiers for request and response tracking.
    \item Error handling for failed publications, invalid payloads, and unavailable subscribers.
\end{itemize}

Zenoh communication must be instrumented with OpenTelemetry spans and Pydantic Logfire logs from the beginning. Each published or consumed message should be observable through both tracing and structured logging.

\section{Alternative Decisions}

One alternative was to route all agent communication through the chat application. This was rejected because the chat application is a user-facing interface, not the dedicated system messaging layer.

Another alternative was to evaluate a different messaging system during Sprint 4. This was rejected because Zenoh had already been retained from prior architectural decisions due to its performance characteristics, topic-based publish-subscribe model, and suitability for future mTLS/PKI-secured transport.

\section{Consequences}

\subsection{Positive}

\begin{itemize}
    \item Agents can communicate asynchronously without direct coupling.
    \item The messaging layer remains independent of the chat application interface.
    \item Topic conventions and message schemas provide a foundation for future agent scaling.
    \item Zenoh interactions can be traced and logged consistently from the beginning.
\end{itemize}

\subsection{Negative}

\begin{itemize}
    \item Topic naming, schema design, and correlation conventions must be documented clearly.
    \item Invalid payloads and unavailable subscribers require explicit error-handling behavior.
    \item Additional integration complexity is introduced during Sprint 4.
\end{itemize}

\section{Scope}

\begin{table}[ht]
\centering
\renewcommand{\arraystretch}{1.25}
\begin{tabular}{L{0.32\textwidth} L{0.58\textwidth}}
\toprule
\textbf{Deliverable} & \textbf{Description} \\
\midrule
Topic conventions & Define Zenoh topics for discovery, commands, responses, status updates, and errors. \\
Message schemas & Define structured payloads for AI-agent-to-AI-agent communication. \\
Publish-subscribe flows & Enable agents to exchange messages asynchronously over Zenoh. \\
Message correlation & Add identifiers for request-response tracking and trace correlation. \\
Error handling & Handle failed publications, invalid payloads, and unavailable subscribers. \\
\bottomrule
\end{tabular}
\end{table}

\newpage

\section*{Statement on AI Use}

Artificial intelligence was used to support the drafting of parts of this document. The generated text was subsequently reviewed, checked, and edited where necessary by me, the author, Abel-Raul Mazilu.

\end{document}